% !TeX document-id = {4a66a092-a797-425f-b588-c0b0f6bcb4d3}
% !TeX TXS-program:compile = txs:///lualatex | txs:///bibtex | txs:///lualatex | txs:///lualatex | txs:///view
% arara: clean: { extensions: ['aux', 'blg', 'bbl', 'log', 'out', 'pdf', 'toc'] }
% arara: lualatex
% arara: bibtex
% arara: lualatex
% arara: lualatex
% arara: clean: { extensions: ['aux', 'blg', 'bbl', 'log', 'out', 'toc'] }

\documentclass{GEM-JNL-FMP}

\usepackage{graphicx}
\usepackage{multicol,multirow}
\usepackage{amsmath,amssymb,amsfonts}
\usepackage{mathrsfs}
\usepackage{amsthm}
\usepackage{rotating}
\usepackage{appendix}
\usepackage[numbers]{natbib}
\usepackage{ifpdf}
\usepackage[T1]{fontenc}
\usepackage{newtxtext}
\usepackage{newtxmath}
\usepackage{textcomp}
\usepackage[colorlinks,allcolors=weblmcolor]{hyperref}

\newtheorem{theorem}{Teorema}[section]
\newtheorem{lemma}[theorem]{Lema}
\theoremstyle{definition}
\newtheorem{remark}[theorem]{Observación}
\newtheorem{example}[theorem]{Ejemplo}
\numberwithin{equation}{section}
\renewcommand{\proofname}{Demostración}
\DeclareMathAlphabet{\mathbb}{U}{msb}{m}{n}
\DeclareMathAlphabet{\mathcal}{OMS}{cmsy}{m}{n}

\artid{1}
\jyear{2025}
\jvol{1}
\jissue{1}

\begin{document}

\begin{Frontmatter}

\title{Resolución Numérica de la ecuación de Lane-Emden}

\author[1]{Matihus Alberth Molina Larios}\orcid{https://orcid.org/0000-0002-9582-9051}
\author[2]{Carlos Alonso Aznarán Laos}\orcid{https://orcid.org/0000-0001-8314-2271}

\authormark{Matihus Alberth Molina Larios \textit{et al}.}

\address[1]{\orgdiv{Facultad de Ciencias Matemáticas}, \orgname{Universidad Nacional Mayor de San Marcos}, \orgaddress{\street{Calle Germán Amézaga 375, Lima}, \postcode{15081}, \country{Perú}};\email{matihus.molina@unmsm.edu.pe}}
\address[2]{\orgdiv{Facultad de Ciencias}, \orgname{Universidad Nacional de Ingeniería}, \orgaddress{\street{Av. Túpac Amaru N$^{\circ}210$ Rímac, Lima}, \postcode{15032}, \country{Perú}};\email{caznaranl@uni.pe}}

\received{4 de diciembre de 2025}
%\revised{10 de diciembre de 2025}
\accepted{12 de diciembre de 2025}

\keywords{Lane-endem, ADM}

\abstract{
    La ecuación de Lane-Emden es un modelo clásico de ecuaciones
    diferenciales ordinarias no lineales con simetría esférica que
    describe la distribución de densidad en el equilibrio en una
    esfera de un gas isotérmico politrópico.
    En este trabajo se analizan sus principales propiedades
    matemáticas, enfatizando en los métodos empleados para estudiar la
    existencia, unicidad y el comportamiento de las soluciones.
    Asimismo, se discuten las implicaciones del índice politrópico
    en la estructura del problema y se presentan resultados que
    evidencian la riqueza analítica del modelo.
}

\end{Frontmatter}

\localtableofcontents

\vspace*{14pt}
\section{Introducción}

La ecuación de Lane-Emden~\cite{lane1870,emden1907} constituye uno de
los modelos más estudiados dentro del ámbito de las ecuaciones
diferenciales ordinarias no lineales con simetría esférica.
Si bien su origen se encuentra en la astrofísica, al modelar la
estructura interna de esferas gaseosas autogravitantes, su estudio
trasciende el contexto físico para convertirse en un problema
matemático de gran interés.
En particular, esta ecuación permite explorar fenómenos de
existencia, unicidad, regularidad y comportamiento asintótico de
soluciones en presencia de singularidades, lo cual la convierte en un
caso paradigmático dentro del análisis de ecuaciones diferenciales
no lineales.

\subsection{Objetivos del trabajo}

El objetivo principal de este trabajo es realizar un estudio matemático detallado de la ecuación de Lane--Emden, destacando sus propiedades estructurales y los métodos utilizados para describir el comportamiento cualitativo de sus soluciones.  
En particular, se busca:
\begin{enumerate}
    \item Analizar la formulación del problema y su interpretación como una ecuación diferencial ordinaria singular.
    \item Estudiar los casos en los que existen soluciones analíticas exactas.
    \item Explorar el comportamiento de las soluciones mediante aproximaciones numéricas y métodos de análisis asintótico.
\end{enumerate}

De este modo, se pretende ofrecer una visión integral del modelo, comprendiendo tanto sus fundamentos teóricos como las herramientas que permiten su estudio riguroso.

\subsection{Estructura del trabajo}

El presente documento se organiza de la siguiente manera:  
En la Sección~2 se desarrolla el marco teórico necesario, abordando conceptos fundamentales sobre ecuaciones diferenciales no lineales y simetría esférica.  
En la Sección~3 se presenta el análisis cualitativo y cuantitativo de la ecuación de Lane--Emden, junto con la discusión de los casos particulares y métodos de resolución.  
Finalmente, en la Sección~4 se exponen las conclusiones generales y posibles líneas de investigación futura relacionadas con el estudio matemático del modelo.

\begin{theorem}[Picard--Lindelöf]
Sea \( D \subset \mathbb{R}^2 \) un dominio abierto y sea \( f : D \to \mathbb{R} \) una función continua.  
Supongamos que existe una constante \( L > 0 \) tal que para todo \( (x, y_1), (x, y_2) \in D \) se cumple la condición de Lipschitz en la segunda variable:
\[
|f(x, y_1) - f(x, y_2)| \leq L |y_1 - y_2|.
\]
Entonces, para todo punto \( (x_0, y_0) \in D \), existe un intervalo \( I = (x_0 - h, x_0 + h) \) y una única función \( y : I \to \mathbb{R} \) de clase \( C^1 \) que satisface el problema de valor inicial
\[
\begin{cases}
y'(x) = f(x, y(x)), \\[6pt]
y(x_0) = y_0.
\end{cases}
\]
En consecuencia, la ecuación diferencial ordinaria posee una \emph{solución local única} que depende continuamente de las condiciones iniciales.
\end{theorem}

La ecuación de Lane--Emden en su forma adimensional se expresa como
\[
\frac{1}{r^{2}} \frac{d}{dr} \left( r^{2} \frac{d\theta}{dr} \right) + \theta^{n} = 0,
\]
donde \( \theta = \theta(r) \) es una función escalar dependiente de la variable radial \( r \geq 0 \), y \( n \) es el \emph{índice politrópico}, que determina el tipo de no linealidad del sistema.  
El problema se complementa con las condiciones de frontera
\[
\theta(0) = 1, \quad \theta'(0) = 0,
\]
que garantizan la regularidad de la solución en el origen.

% TODO: Mostrar algunas soluciones exactas dependen del índice.
Existen una inmensidad de métodos para resolver estas EDP's,mencionaremos a las más importantes.
\section{Método 0}
\textbf{Daremos una explicación breve del metodo de Adominam y como es que funciona.\cite{adomian1995}.}

El método de Adomiam es muy usado para aproximar ecuaciones diferenciales no lineales, mediante sumatorias y aproximaciones.Es decir, la solución $y(r)$ como una serie:
$$y(r)=\sum_{k=0}^{\infty}y_k(r)$$
y descomponer la parte no lineal $N(y)$ en una serie de polinomios de Adomiam:
$$N(y(r))=\sum_{k=0}^{\infty}A_k(r)$$
En el paper de Adomian \cite{adomian1995} esta el método en general,procederemos a citar el método aplicado para la ecuación de lane-endem.

\section{Aplicación del método de descomposición (modificado) al problema de Lane--Emden}

Consideremos la forma general (politrópica) de la ecuación de Lane--Emden escrita para la función \(T=T(r)\):
\[
\frac{1}{r^{2}}\frac{d}{dr}\!\left(r^{2}\frac{dT}{dr}\right) + T^{m} = 0,\qquad r\ge 0,
\]
con las condiciones iniciales regulares en el origen
\[
T(0)=1,\qquad T'(0)=0,
\]
donde \(m\in\mathbb{N}\) parametriza la no linealidad.

Se supone una solución en serie
\[
T(r)=\sum_{k=0}^{\infty} T_{k}(r),
\]
y se descompone la parte no lineal en polinomios de Adomian \(A_k\):
\[
T(r)^{m} = \sum_{k=0}^{\infty} A_{k}(r).
\]

Con estas condiciones iniciales concluimos que $a_0=1 ,a_1=0$.

Y que $r^{-2}\frac{d}{dr}r^2\frac{d}{dr}T=\sum_{n=0}^{\infty}(n+2)(n+3)a_{n+2}r^n$
Definimos al operador lineal $L=r^{-2}\frac{d}{dr}r^2\frac{d}{dr}$ y $N(T)=\lambda T^m$, donde $LT+NT=0$.Luego:
$$\sum_{n=0}^{\infty}(n+2)(n+3)a_{n+2}r^n+\lambda\sum_{n=0}^{\infty}A_n(r)r^n=0$$.
$$(n+2)(n+3)a_{n+2}r^n+\lambda A_n(r)r^n=0$$

Con ello ya podemos encontrar los terminos de $a$ y de $A$.










Utilizaremos el metodo propuesto por wazwaz \cite{wazwaz} .
En este paper generaliza el metodo de Adominam que ya es muy usado en la física,además dirrección el problema tratando con el punto singular x=0.
\section{Método 1}
Aquí nos centraremos en el metodo propuesto por el colombiano Cardenas.
En este aborda para $f(x)=e^{+-x}$ y $f(x)=(x^2-C)^{\frac{3}{2}}$.
Este método se basa en la misma lógica que la Transformadas de Laplace.Definimos. el operador:
$$Y(k)=\frac{1}{k!}[\frac{d^ky(x)}{dx^k}]_{x=x_0}$$
$y(x)=\sum_{k=0}^{\infty}Y(k)x^k$,junto con otros teoremas.
\section{Análisis de convergencia}
Encontramos un Analisis de convergencia en el trabajo de Cardénas \cite{alzate2014}
\begin{theorem}[Convergencia del método de descomposición]
Sea $\mathcal{T}$ un operador no lineal definido en un espacio de Banach y sea $y(x)$ la solución exacta de la ecuación
\begin{equation}\label{eq:42}
\mathcal{T}(y) = y.
\end{equation}
La serie solución
\begin{equation}\label{eq:14}
y(x) = \sum_{k=0}^{\infty} y_k(x)
\end{equation}
converge a $y(x)$, si existe una constante $c$ tal que $0 \leq c < 1$ y
\[
\|y_{k+1}\| \leq c\,\|y_k\|, \qquad \forall\, k \in \mathbb{Z}_{\ge 0}.
\]
\end{theorem}

\begin{proof}
Demostremos que la sucesión $\{T_n\}_{n=0}^{\infty}$ es de Cauchy en el espacio de Banach considerado.  
Por hipótesis se tiene:
\[
\|T_{n+1} - T_n\| = \|y_{n+1}\| \leq c\,\|y_n\| \leq c^2\,\|y_{n-1}\| \leq \dots \leq c^{n+1}\|y_0\|.
\]
Entonces, para cualesquiera $m,n \in \mathbb{N}$ con $m > n$, se cumple:
\[
\begin{aligned}
\|T_m - T_n\| 
&= \| (T_m - T_{m-1}) + (T_{m-1} - T_{m-2}) + \cdots + (T_{n+1} - T_n) \| \\[4pt]
&\le \|T_m - T_{m-1}\| + \|T_{m-1} - T_{m-2}\| + \cdots + \|T_{n+1} - T_n\| \\[4pt]
&\le c^{m}\|y_0\| + c^{m-1}\|y_0\| + \cdots + c^{n+1}\|y_0\| \\[4pt]
&= c^{n+1}\|y_0\| \left(1 + c + c^2 + \cdots + c^{m-n-1}\right) \\[4pt]
&= c^{n+1}\|y_0\|\,\frac{1 - c^{m-n}}{1-c} \\[4pt]
&< \frac{c^{n+1}}{1-c}\,\|y_0\|.
\end{aligned}
\]
Dado que $0 \le c < 1$, el término del lado derecho tiende a cero cuando $n \to \infty$.  
Por tanto, $\{T_n\}$ es una sucesión de Cauchy, y como el espacio es completo, existe $T \in X$ tal que
\[
\lim_{n\to\infty} T_n = T.
\]
Así, la serie
\[
\sum_{n=0}^{\infty} y_n = T
\]
converge.

Además, dado que la ecuación \eqref{eq:42} equivale a resolver $\mathcal{T}(T) = T$, si el operador $\mathcal{T}$ es continuo, se tiene
\[
\lim_{n\to\infty} T_{n+1} = \lim_{n\to\infty} \mathcal{T}(T_n) = \mathcal{T}\!\left( \lim_{n\to\infty} T_n \right) = \mathcal{T}(T).
\]
Por lo tanto, $T$ es una solución de
\[
y(x) = \mathcal{T}(y(x)),
\]
y el teorema queda demostrado.
\end{proof}

\section{Ejemplos numéricos}

\section{Conclusiones}



\bibliographystyle{plainnat}
\bibliography{references}


\end{document}
